\section{Matrices}
	
	We have developed a good theory of vector spaces and linear maps so far. But in order to do computations in practice, we should identify a concrete realization of vectors and linear maps. We will focus on finite-dimensional vector spaces. 

	\begin{definition}[Column Vector]
		Let $V$ be a finite-dimensional vector space over $\mathbb{F}$ and $\{v_1, \ldots, v_n\}$ be a basis of $V$. Then, any vector $v \in V$ can be written as the linear combination 
		\begin{equation}
		  v = a_1 v_1 + \ldots + a_n v_n
		\end{equation}
		Since $V$ is isomorphic to $\mathbb{F}^n$, we can associate such $v$ with the \textbf{column vector}
		\begin{equation}
			v \mapsto \begin{pmatrix} a_1 \\ \vdots \\ a_n \end{pmatrix}
		\end{equation}
	\end{definition} 

	Therefore, as long as we choose a basis of $V$, we can basically treat the vectors as $\mathbb{F}^n$. What about linear maps from $V$ to $W$? First, let's put a basis on each of them $\{v_1, \ldots, v_n\}$ and $\{w_1, \ldots, w_m\}$ and identify the isomorphisms: 
	\begin{align}
		\iota_V : V \to \mathbb{F}^n, \qquad v = \sum_{i=1}^n a_i v_i \mapsto \begin{pmatrix} a_1 \\ \vdots \\ a_n \end{pmatrix} \in \mathbb{F}^n \\
		\iota_W : W \to \mathbb{F}^m, \qquad w = \sum_{j=1}^m b_j w_j \mapsto \begin{pmatrix} b_1 \\ \vdots \\ b_m \end{pmatrix} \in \mathbb{F}^m
	\end{align}

	Let's first identify how $A$ should map the $n$ basis vectors to. That is, just choose any $n$ vectors $u_1, \ldots, u_n \in W$ and set them to be 
	\begin{equation}
	  u_1 = A(v_1), u_2 = A(v_2), \ldots, u_n = A(v_n) \in W
	\end{equation}
	By linearity of $A$, this completely determines the map. This is because for every $v \in V$, 
	\begin{equation}
		A(v) = A \bigg( \sum_{i=1}^n a_i v_i \bigg) = \sum_{i=1}^n a_i A(v_i) = \sum_{i=1}^n a_i u_i
	\end{equation}

	But since each $u_i$ is a vector in $W$ and $W$ has a basis, we can expand it with respect to the basis of $W$. 
	\begin{equation}
		A(v_i) = \sum_{j=1}^m a_{ji} w_j, \qquad i = 1, \ldots, n
	\end{equation} 

	Therefore, we have identified a \textit{new} map $M: \mathbb{F}^n \to \mathbb{F}^m$ such that 
	\begin{equation}
		A = \iota_W^{-1} \circ M \circ \iota_V
	\end{equation} 
	So basically, any linear map between finite-dimensional vector spaces can be treated as a linear map from $\mathbb{F}^n$ to $\mathbb{F}^m$. For convenience, we encode this linear map as a rectangular array of numbers. 

	\begin{definition}[Matrix]
		Let $V, W$ be finite dimensional vector spaces over $\mathbb{F}$, with bases $B_V = \{v_1, \ldots, v_n\}$ and $B_U = \{u_1, \ldots, u_m\}$, respectively. Let $A: V \to W$ be a linear map. Then, the \textbf{matrix realization} of linear map $A$ with respect to bases $B_V$ and $B_U$\footnote{Note that $M_A$ is really dependent on the bases, but generally, we omit the basis if it is clear from context.} is 
		\begin{equation}
			M_A \coloneqq \begin{pmatrix}
				a_{11} & a_{12} & \ldots & a_{1n} \\ 
				a_{21} & a_{22} & \ldots & a_{2n} \\ 
				\vdots & \vdots & \ddots & \vdots \\ 
				a_{m1} & a_{m2} & \ldots & a_{mn} 
			\end{pmatrix} \in \mathbb{F}^{m \times n}
		\end{equation} 
		where the elements satisfy 
		\begin{equation}
			A(v_i) = \sum_{j=1}^m a_{ji} w_j \in W
		\end{equation}
	\end{definition}

	Therefore, the action of linear map $A$ on vector $v \in V$ can be realized as the \textit{left matrix multiplication} of $M_A$ on vector $\iota_V (v) \in \mathbb{F}^n$. It is important to note that the matrix is \textit{not} $A$ in of itself. In the most rigorous sense, the matrix $A$ is really just equal to the composition of mappings $\iota_W^{-1} A \iota_V$, but for simplicity it is just written as $A$. It is just one representation of a linear map given the two bases of the domain and codomain. As soon as one writes down a matrix to represent a linear map, they are automatically assuming some choice of basis given by $\iota_V$ and $\iota_W$. 

	Note that the way we represent vectors and linear transformation has all been arbitrarily chosen. There is nothing innate about the way we express these transformation as matrix multiplication. It is dependent on our \textit{initial choice} to represent linear mappings as \textit{left} matrix multiplication and to represent all vectors as column vectors. This leads us to write dual vectors $l$ as \textit{row vectors}, since we can interpret their action on a vector $v$ as left matrix multiplication as well. 	
	\begin{equation}
	  l(v) = lv
	\end{equation}

	\begin{definition}[Row Vector]
		Let $V$ be a finite-dimensional vector space over $\mathbb{F}$ and $V^\ast$ be the dual space of $V$ with basis $\{l_1, l_2, \ldots, l_n\}$. Then, any vector $l \in V^\ast$ can be written as the linear combination 
		\begin{equation}
		  v = a_1 l_1 + \ldots + a_n l_n
		\end{equation}
		Since $V^\ast$ is isomorphic to $\mathbb{F}^n$, we can assocaite such $l \in V^\ast$ with the \textbf{row vector} 
		\begin{equation}
			l \mapsto \begin{pmatrix} a_1 & \ldots & a_n \end{pmatrix}
		\end{equation}
	\end{definition}

  Numerically, when computing the product two $n \times n$ matrices $A$ and $B$ to another $n \times n$ matrix $C$, since each entry of $C$ is the product of a row of $A$ with a column of $B$, and since $C$ has $n^2$ entries, we need $n^3$ scalar multiplications to compute (as well as $n^3 - n^2$ additions). In order words, the computing efficiency of the algorithm is at $O(n^3)$. However, there are faster algorithms than this. This is algorithm is known as the \textbf{Strassen Algorithm} (however, there do exist faster algorithms). 

  \begin{algo}[Strassen Algorithm]
    Let $A, B$ be $2 \times 2$ matrices such that $AB = C$. That is, component-wise,
    \begin{equation}
      \begin{pmatrix}
      a_{11} & a_{12} \\ a_{21} & a_{22}
      \end{pmatrix} \begin{pmatrix}
      b_{11} & b_{12} \\ b_{21} & b_{22}
      \end{pmatrix}
       = \begin{pmatrix}
        c_{11} & c_{12} \\ c_{21} & c_{22}
       \end{pmatrix}
    \end{equation}
    where for $i, j = 1, 2$, 
    \begin{equation}
      c_{ij} = a_{i1} b_{1j} + a_{i2} + b_{2j}
    \end{equation}
    Then, let us define 
    \begin{align*}
      P_1 &= (a_{11} + a_{22}) (b_{11} + b_{22}) \\
      P_2 &= (a_{21} + a_{22}) b_{11} \\
      P_3 &= a_{11} (b_{12} - b_{22}) \\
      P_4 &= a_{22} (b_{21} - b_{11} \\
      P_5 &= (a_{11} + a_{12}) b_{22} \\
      P_6 &= (a_{21} - a_{11}) (b_{11} + b_{12}) \\
      P_7 &= (a_{12} - a_{22}) (b_{21} + b_{22}) 
    \end{align*}
    Then, the theorem states that we the entries of $C$ are 
    \begin{align*}
      c_{11} &= P_1 + P_4 - P_5 + P_7 \\
      c_{12} &= P_3 + P_5 \\
      c_{21} &= P_2 + P_4 \\
      c_{22} &= P_1 + P_3 - P_2 + P_6
    \end{align*}

		This algorithm for multiplying $2\times 2$ matrices requires $7$ scalar multiplications, while regular multiplication requires $8$. Using block multiplication, we can use this algorithm to calculate any matrix of order $2^k$. That is, to calculate $2^k \times 2^k$ matrices, we have to perform seven multiplications of blocks of size $2^{k-1} \times 2^{k-1}$, and doing this recursively, it reduces it down to 
		\begin{equation}
			7^k = 2^{k \log_2{7}} = n^{\log_2{7}}
		\end{equation}
		where $n$ is the order of the matrices being multiplied. 

		Additionally, the number of scalar additions or subtractions needed is bounded by 
		\begin{equation}
			6 \times 7^k = 6 \times 2^{k \log_2{7}} = 6 n^{\log_2{7}}
		\end{equation}
		Since $\log_2{7} \approx 2.807 < 3$, this algorithm does indeed have more computational efficiency.\footnote{It is conjectured that for any positive number $\varepsilon$, there is an algorithm that computes the product of two $n \times n$ matrices with computational efficiency of $O(n^{2 + \varepsilon})$. } Note that matrices whose order is not a power of $2$ can be turned into one by adjoining a suitable number of $1$s on the diagonal.
  \end{algo}

	\begin{definition}[Column Space]
		The \textbf{column space} of an $m \times n$ matrix 
		\begin{equation}
			A = \begin{pmatrix} | & \ldots & | \\ 
				a_1 & \ldots & a_n \\ 
				| & \ldots & |
			\end{pmatrix}
		\end{equation}
		is called $C(A) \coloneqq \mathrm{span}\{a_1, \ldots, a_n\}$. 
	\end{definition}

	\begin{definition}[Null Space]
		The \textbf{null space} of a matrix $\mathbb{F}^{m \times n}$ is the set 
		\begin{equation}
			N(A) \coloneqq \{ v \in \mathbb{F}^n \colon Av = 0 \}
		\end{equation}
	\end{definition}

\subsection{Algebra of Matrices} 

	We hope that this set of all $m \times n$ matrices, denoted $\mathbb{F}^{m \times n}$, satisfies the same properties of the set of linear maps $\mathcal{L}(V, W)$. Indeed, it has a vector space structure isomorphic to the vector space of linear maps. 


	We want to add the additional vector space and algebra structure on these sets of matrices. 

	\begin{theorem}[Vector Space of Matrices]
		The set of all $m \times n$ matrices over $\mathbb{F}$, denoted $\mathbb{F}^{m \times n}$, is a vector space over $\mathbb{F}$, with the following operations. 
		\begin{enumerate}
			\item \textit{Addition}. Given $A, B \in \mathbb{R}^{m \times n}$, 
				\begin{equation}
					(A + B)_{ij} = A_{ij} + B_{ij}
				\end{equation}
			\item \textit{Scalar Multiplication}. Given $A \in \mathbb{R}^{m \times n}$ and $c \in \mathbb{F}$, 
				\begin{equation}
					(cA)_{ij} = c A_{ij}
				\end{equation}
		\end{enumerate}
		Furthermore, $\mathbb{F}^{m \times n} \simeq \mathcal{L}(\mathbb{F}^n, \mathbb{F}^m)$ as vector spaces. 
	\end{theorem}

	Furthermore, it is the case that the composition operation in the algebra of linear operators is realized as the operation of matrix multiplication. These are two distinct operations that are related only through the basis mappings $i$ and $j$. 

	\begin{definition}[Matrix Multiplication]
		Given two matrices $A \in \mathbb{F}^{m \times n}$ and $B \in \mathbb{F}^{n \times p}$, the \textbf{matrix multiplication} of $A$ and $B$ is defined as the matrix $AB$ where 
		\begin{equation}
			(AB)_{ij} = \sum_{k=1}^n A_{ik} B_{kj}
		\end{equation}
	\end{definition}

	\begin{lemma}[Matrix Multiplication as Realization of Composition]
	  Let $U, V, W$ be finite-dimensional vector spaces each with their respective bases and $T: U \to V, S: V \to W$ be linear maps. Then, 
		\begin{equation}
			M_{ST} = M_S M_T
		\end{equation}
		That is, the matrix realization of the composition $ST$ is the matrix product of the realizations $S$ and $T$. 
	\end{lemma}
	\begin{proof}
	  
	\end{proof}

	\begin{theorem}[Algebra of Matrices]
		The set of all $n \times n$ matrices over field $\mathbb{F}$, denoted $\mathbb{F}^{n \times n}$, is a noncommutative associate algebra isomorphic to $\mathcal{L}(\mathbb{F}^n)$. 
  \end{theorem}

	\begin{example}[Trigonometric Identities from Rotations]
    Let $\alpha: \mathbb{R}^2 \longrightarrow \mathbb{R}^2$ be the linear transformation of the counterclockwise rotation by $\theta$ and $\beta: \mathbb{R}^2 \longrightarrow \mathbb{R}^2$ be the counterclockwise rotation of $\phi$. Then the matrix representation of $\alpha \circ \beta$ is 
    \begin{align}
        & \begin{pmatrix}
      \cos{\theta} & - \sin{\theta} \\
      \sin{\theta} & \cos{\theta}
      \end{pmatrix} \begin{pmatrix}
      \cos{\phi} & - \sin{\phi} \\
      \sin{\phi} & \cos{\phi} 
      \end{pmatrix} \\
       & = \begin{pmatrix}
      \cos{\theta} \cos{\phi} - \sin{\theta} \sin{\phi} & - \sin{\phi} \cos{\theta} - \cos{\phi} \sin{\theta} \\
      \sin{\theta} \cos{\phi} + \cos{\theta} \sin{\phi} & - \sin{\theta} \sin{\phi} + \cos{\theta} \cos{\phi}
      \end{pmatrix}
    \end{align}
    But the counterclockwise rotation by $\theta$ and then $\phi$ is really just a counterclockwise rotation by $\theta + \phi$, which has the matrix representation
    \begin{equation}
      \begin{pmatrix}
      \cos{(\theta + \phi)} & - \sin{(\theta + \phi)} \\
      \sin{(\theta + \phi)} & \cos{(\theta + \phi)}
      \end{pmatrix}
    \end{equation}
    Since both matrices must be equivalent, this produces the trigonometric identities for angle addition.
    \begin{align*}
      \sin{(\theta + \phi)} = \sin{\theta} \cos{\phi} + \cos{\theta} \sin{\phi} \\
      \cos{(\theta + \phi)} = \cos{\theta} \cos{\phi} - \sin{\theta} \sin{\phi}
    \end{align*}
  \end{example}

	Unsurprisingly, we can define a Euclidean norm on the vector space of matrices.

	\begin{definition}[Frobenius Norm]
	  The \textbf{Frobenius norm} of a matrix $A$ is defined as 
		\begin{equation}
			\|A\|_F = \sqrt{\sum_{i, j} a_{ij}^2}
		\end{equation}
	\end{definition}

	A natural question to ask is how this is related to the operator norm for general linear maps? It turns out that it is an upper bound. 

	\begin{theorem}[Frobenius Norm vs Operator Norm]
	  For any matrix $A$, 
		\begin{equation}
		  \|A\| \leq \|A\|_F
		\end{equation}
	\end{theorem}
	\begin{proof}
	  
	\end{proof}

\subsection{Row Reduction}

	Matrices are generally unruly to work with since they require a lot of computation. Therefore, we would like to transform them into some ``nicer'' matrices. Let's first define some nice matrices. 

	\begin{definition}[Diagonal Matrix]
		A square matrix $A$ is \textbf{diagonal} if $a_{i j} = 0$ for $i \neq j$. 
	\end{definition}

	\begin{definition}[Triangular Matrices]
    A matrix $A$ is called 
		\begin{enumerate}
		  \item \textbf{lower triangular} if $a_{i j} = 0$ for $i < j$. 
			\item \textbf{upper triangular} if $a_{i j} = 0$ for $i > j$. 
		\end{enumerate}
  \end{definition}

  \begin{definition}[Pivot]
		Let $A \in \mathbb{F}^{m \times n}$. Then, 
		\begin{enumerate}
		  \item The \textbf{pivot} of a row $(a_1, a_2, \cdots, a_n)$ is its first nonzero element. 
			\item If this element is $a_k$, then $k$ is the \textbf{index} of the pivot. 
			\item The rest of the columns that don't contain a pivot are called \textbf{free variables}. 
		\end{enumerate}
  \end{definition}

  \begin{definition}[Row Echelon Form]
    A matrix is in \textbf{Echelon form}, or \textbf{row Echelon form}, if 
    \begin{enumerate}
      \item the indices of the pivots of its nonzero rows form a strictly increasing sequence, like steps
      \item zero rows, if they exist, are at the bottom
    \end{enumerate}

		\begin{figure}[H]
			\centering 
			\begin{equation}
				\begin{pmatrix}
					a_{1 j_1} & * & \ldots & \ldots & * \\
					0 & a_{2 j_2} & * & \ldots & * \\
					0& 0& \ddots & \ldots & * \\
					0& 0& 0& a_{r j_r} & \vdots \\
					0 & 0 & \ldots & 0 & 0
				\end{pmatrix}
			\end{equation}
			\caption{A matrix in Echelon form is in form where $*$'s represent arbitrary numbers, $a_{i j_i}$'s are nonzero (with indices $j_i$, and the entries to the left of below them are $0$. Property $(i)$ also states that $j_1 < j_2 < \cdots < j_r$. } 
			\label{fig:echelon}
		\end{figure}
  \end{definition}

	\begin{definition}[Reduced Row Echelon Form]
    A matrix is in \textbf{reduced row echelon form}, denoted rref$(A)$, if
    \begin{enumerate}
      \item it is in row echelon form
      \item the pivots are all equal to 1
      \item each column containing a pivot has zeros in all other entries
    \end{enumerate}
  \end{definition}

  \begin{definition}[Elementary Row Transformation]
    An \textbf{elementary row transformation} of a matrix $A$ is one of the following three types of transformations
    \begin{enumerate}
      \item \textit{Scalar Multiplication}. Multiplying a row by a nonzero scalar: $A_i \gets c A_i$. 
      \item \textit{Swap Rows}. Interchanging two rows: $A_i \gets A_j, A_j \gets A_i$. 
      \item \textit{Scaled Addition} Adding a row multiplied by a scalar to another row: $A_i \gets A_i + c A_j$. 
    \end{enumerate}
  \end{definition} 

	\begin{example}[Elementary Row Transformations]
	  We present an example of each of the three transformations. Let
		\[
			A =
			\begin{pmatrix}
				1 & 2 & 3 \\
				4 & 5 & 6 \\
				7 & 8 & 9
			\end{pmatrix}.
		\]
		\begin{enumerate}
		  \item \textit{Scalar Multiplication}. If we multiply the second row by $2$, then
			\[
				A_2 \gets 2 A_2
			\]
			gives
			\[
				\begin{pmatrix}
					1 & 2 & 3 \\
					4 & 5 & 6 \\
					7 & 8 & 9
				\end{pmatrix}
				\longrightarrow
				\begin{pmatrix}
					1 & 2 & 3 \\
					8 & 10 & 12 \\
					7 & 8 & 9
				\end{pmatrix}.
			\]

			\item \textit{Swap Rows}. If we interchange the first and third rows, then
			\[
				A_1 \leftrightarrow A_3
			\]
			gives
			\[
				\begin{pmatrix}
					1 & 2 & 3 \\
					4 & 5 & 6 \\
					7 & 8 & 9
				\end{pmatrix}
				\longrightarrow
				\begin{pmatrix}
					7 & 8 & 9 \\
					4 & 5 & 6 \\
					1 & 2 & 3
				\end{pmatrix}.
			\]

			\item \textit{Scaled Addition}. If we add $-4$ times the first row to the second row, then
			\[
				A_2 \gets A_2 - 4 A_1
			\]
			gives
			\[
				\begin{pmatrix}
					1 & 2 & 3 \\
					4 & 5 & 6 \\
					7 & 8 & 9
				\end{pmatrix}
				\longrightarrow
				\begin{pmatrix}
					1 & 2 & 3 \\
					0 & -3 & -6 \\
					7 & 8 & 9
				\end{pmatrix}.
			\]
		\end{enumerate}
	\end{example}

	\begin{lemma}[Elementary Row Transformations as Matrix Multiplication]
		The elementary transformations on a $m \times n$ matrix $A$ is equivalent to left matrix multiplication by the following $m \times m$ matrices.\footnote{Due to the following difficulty in presenting these matrices in a general form, we present them in the specific $4 \times 4$ case and hope that the reader can extrapolate this process to general matrices. }

		\begin{enumerate}
			\item Adding row $i$ multiplied by scalar $\alpha$ to row $j$ (where $j > i$) is denoted $E^1_{\alpha \times i + j}$. The matrix is the identity matrix with $\alpha$ in the $(j, i)$ position. 
			\begin{equation}
				E^1_{2 \times 1 + 2} = \begin{pmatrix}
				1&0&0&0 \\ 2&1&0&0 \\ 0&0&1&0 \\ 0&0&0&1
				\end{pmatrix}, \;\; E^1_{-3 \times 2 + 4} = \begin{pmatrix}
				1&0&0&0 \\ 0&1&0&0 \\ 0&0&1&0 \\ 0&-3&0&1
				\end{pmatrix}
			\end{equation}
			\item Interchanging the $i$th and $j$th row is denoted by matrix $E^2_{i j}$. Note that these are permutation matrices, or more speficially, transpositions. 
			\begin{equation}
				E^2_{2 3} = \begin{pmatrix}
				1&0&0&0 \\ 0&0&1&0 \\ 0&1&0&0 \\ 0&0&0&1
				\end{pmatrix}, \; \; E^2_{2 4} = \begin{pmatrix}
				1&0&0&0 \\ 0&0&0&1 \\ 0&0&1&0 \\ 0&1&0&0
				\end{pmatrix}
			\end{equation}

			\item Multiplying the $i$th row by a scalar $\alpha$ is denoted by matrix $E^3_{\alpha \times i}$. 
			\begin{equation}
				E^3_{3 \times 3} = \begin{pmatrix}
				1&0&0&0 \\ 0&1&0&0 \\ 0&0&3&0 \\ 0&0&0&1
				\end{pmatrix}, \;\; E^3_{7 \times 1} = \begin{pmatrix}
				7&0&0&0 \\ 0&1&0&0 \\ 0&0&1&0 \\ 0&0&0&1
				\end{pmatrix}
			\end{equation}
		\end{enumerate}
	\end{lemma}
	\begin{proof}
	  Trivial. 
	\end{proof}

	\begin{lemma}[Inverse of Elementary Row Transformations]
    Each elementary matrix is invertible and their inverses are also elementary matrices. More specifically, 
    \begin{enumerate}
      \item $(E^1_{\alpha \times i + j})^{-1} = E^1_{-\alpha \times i + j}$ (same matrix but $\alpha$ changed to -$\alpha$)
      \item $(E^2_{i j})^{-1} = E^2_{i j}$ (same matrix) 
      \item $(E^3_{\alpha \times i})^{-1} = E^{3}_{(1/\alpha) \times i}$ (same matrix but $\alpha$ changed to $1 / \alpha$)
    \end{enumerate}
  \end{lemma}

  Elementary column operations of are equivalent to right multiplication of matrices. However, since we defined the action of a linear map as left matrix multiplication, we don't focus on column operations. 

	\begin{theorem}[Gauss Jordan Elimination]
    Every matrix can be reduced to row echelon form and reduced row echelon form by elementary row transformations. 
  \end{theorem}
  \begin{proof}
    The relevant algorithm used will not be shown here, but we will mention that this procedure is called \textbf{Gauss Elimination}, or \textbf{row reduction}. Then, we perform the algorithm known as \textbf{back substitution}, where we start with the bottom row and use elementary operations to cancel out terms in upper rows. 
  \end{proof}

  Under the column space interpretation, Gaussian Elimination is really just an algorithm that performs a change of basis in steps. Each elementary operation simultaneously changes all of the vectors of the column space in such a way that eventually, this set of vectors will be "nice-looking" with a lot of zero entries. 

  The computational efficiency of Gauss Elimination is well known. Solving a system of $n$ equations with $n$ variables with this algorithm requires approximately $2 n^3 / 3$ operations, meaning that it has arithmetic complexity of $O(n^3)$. However, for matrices of large order, multiple problems can occur. The algorithm generally does not have memory problems if the field is finite or if the coefficients are floating-point numbers. However, if the coefficients are integers or rational numbers, the intermediate entries of the algorithm can grow exponentially large, so bit complexity is exponential. However, there is a variant of Gaussian elimination, called the Bareiss algorithm, that avoids this problem, but has bit complexity of $O(n^5)$. Another problem is numerical instability, caused by the possibility of dividing by numbers very close to $0$. Any such number would have its existing error amplified. Gaussian elimination algorithm is generally known to be stable for positive-definite matrices. 

  \begin{definition}[Rank]
    The number of pivots in ref$(A)$ is called the \textbf{rank} of $A$, denoted $\rank(A)$. 
  \end{definition}

	\begin{theorem}[Rank as Dimension of Image]
		For any $A \in \mathbb{F}^{m \times n}$, it is true that 
		\begin{equation}
			\rank(A) = \dim(\im(A)) = \dim(C(A)) \leq \min\{m, n\}
		\end{equation}
    There are $k$ free variables in $A$ if and only if $\dim(N(A)) = k$. 
  \end{theorem}
  \begin{proof}
    By definition, the number of pivots cannot exceed the number of variables nor can it exceed the number of equations. Let $A$ be a $m \times n$ matrix over $\mathbb{F}$. Then, let $\rank(A) = k$, which implies that there are $n-k$ free variables $\implies \dim(N(A)) = n - k$. By rank nullity, 
    \begin{equation}
      \dim(\im(A)) = n - \dim(N(A)) = n - (n -k) = k = \rank(A)
    \end{equation}

    For the free variables part, we do not give a rigorous proof but we outline one. Each free variable corresponds to a free vector in the row echelon form of $A$ that are all linearly independent. Since the span of these free vectors is equal to $N(A)$, the $k$ vectors form a basis of $A$ $\implies$ by definition, $\dim(N(A)) = k$.
  \end{proof}

  This theorem establishes the consistency in definition between the rank of an abstract mapping mentioned in chapter 2 and the rank of its matrix representation. We can in fact establish strong claims on top of this. When we have a square matrix, we have additional properties. 

	\begin{definition}[Nonsingular Matrix]
		$A \in \mathbb{F}^{n \times n}$ is called 
		\begin{enumerate}
		  \item \textbf{nonsingular} if $\rank(A) = n$, and 
			\item \textbf{singular} if $\rank(A) < n$. 
		\end{enumerate}
  \end{definition}

	\begin{theorem}[Invertibility]
    $n \times n$ matrix $A$ is invertible if and only if it is nonsingular. 
  \end{theorem}
  \begin{proof}
    $A$ is nonsingular $\iff A x = b$ will always have a unique solution $\iff$ $A$ is an isomorphism from $\mathbb{F}^n$ to itself $\iff$ by definition, $A$ is invertible. 
  \end{proof}

\subsection{Change of Basis}

	So far, we have interpreted linear automorphisms as a map $A: V \to V$ mapping $v$ to $Av$. This interpretation of $A$ as an \textit{active transformation} is the most straightforward. But given a basis $\{v_1, \ldots, v_n\}$, we can see that $A$ is not acting on the coefficients, but the \textit{on the basis vectors themselves}. 
	\begin{equation}
	  Av = A (a_1 v_1 + \ldots + a_n v_n) = a_1 Av_1 + \ldots + a_n A v_n
	\end{equation}

	Therefore, we can interpret the mapping $A$ as a mapping of the basis vectors only, while keeping the coefficients fixed. That is, $Av$ merely represents $v$ with respect to another set of basis vectors $\{A v_1, \ldots, A v_n\}$. If we interpret map as this, then this is called a \textit{passive transformation}. We can interpret $A$ as both an active or a passive transformation. 

  Suppose $\{u_1, u_2, \ldots, u_n\}$ is a basis for vector space $V$ and $\{v_1, v_2, \ldots, v_n\}$ is another basis for $V$. Then, for any vector $v \in V$, we can write 
	\begin{align}
		v & = a_1 u_1 + \ldots + a_n u_n \\ 
			& = b_1 v_1 + \ldots + b_n v_n
	\end{align}
	The realization of $v$ with respect to $\{a_i\}$ will be $a = (a_1, \ldots, a_n)^T$ and the realization of $v$ with respect to $\{b_i\}$ will be $b = (b_1, \ldots, b_n)^T$. Therefore, we want to find some matrix $S$ such that 
	\begin{equation}
		b = Sa \iff \begin{pmatrix} b_1 \\ \vdots \\ b_n \end{pmatrix} = \begin{pmatrix} s_{11} & \ldots & s_{1n} \\  \vdots & \ddots & \vdots \\ s_{n1} & \ldots & s_{nn} \end{pmatrix} \begin{pmatrix} a_1 \\ \vdots \\ a_n \end{pmatrix} 
	\end{equation}
	This matrix certainly exists, and to find matrix $S$ that transforms from the $u_i$'s to the $v_j$'s, we must start off by representing $v = \sum_i a_i u_i$, and trying to present it into a form like $\sum_j b_j v_j$. To do this, we work backwards by writing each term of the \textit{old basis} as linear combinations of the \textit{new basis}. 
	\begin{equation}
		u_i = \sum_{j} s_{ji} v_j
	\end{equation}

	Once we have figured out all values $s_{ji}$, we can rewrite $v$ from the old basis to the new basis. 
	\begin{align}
		v & = \sum_{i} a_i u_i \\ 
			& = \sum_{i} a_i \sum_j s_{ji} v_j \\ 
			& = \sum_{i, j} a_i s_{ji} v_j \\ 
			& = \sum_{j} \bigg( \sum_i a_i s_{ji} \bigg) v_j \\ 
			& = \sum_j b_j v_j
	\end{align} 
	Therefore, we find that 
	\begin{equation}
		b_j = \sum_i s_{ji} a_i 
	\end{equation} 

	\begin{definition}[Change of Basis Matrix]
		For all $v \in V$, the matrix $S$ thaat transforms the coefficients $(a_1, \ldots, a_n)^T$ of $v$ with respect to basis $\{u_1, \ldots, u_n\}$ to the coefficients $(b_1, \ldots, b_n)^T$ of $v$ with respect to basis $\{v_1, \ldots, v_n\}$ is called a \textbf{change-of-basis matrix}. 
	\end{definition}

	\begin{example}[Computation of a Change-of-Basis Matrix]
	  Let 
		\begin{equation}
			\{u_1, u_2, u_3 \} = \bigg\{ \begin{pmatrix} 1 \\ 0 \\ 0 \end{pmatrix}, \begin{pmatrix} 0 \\ 1 \\ 0 \end{pmatrix}, \begin{pmatrix} 0 \\ 0 \\ 1 \end{pmatrix} \bigg\}, \qquad \{v_1, v_2, v_3 \} = \bigg\{ \begin{pmatrix} 1 \\ 1 \\ 0 \end{pmatrix}, \begin{pmatrix} 0 \\ 1 \\ 1 \end{pmatrix}, \begin{pmatrix} 1 \\ 0 \\ 1 \end{pmatrix} \bigg\}
		\end{equation}
		by two bases of $\mathbb{R}^3$. Then, we see that 
		\begin{align}
			u_1 & = \frac{1}{2} v_1 - \frac{1}{2} v_2 + \frac{1}{2} v_3 \\ 
			u_2 & = \frac{1}{2} v_1 + \frac{1}{2} v_2 - \frac{1}{2} v_3 \\ 
			u_3 & = -\frac{1}{2} v_1 + \frac{1}{2} v_2 + \frac{1}{2} v_3 
		\end{align}
		Therefore, our change-of-basis matrix should be 
		\begin{equation}
			S = \begin{pmatrix}
				\frac{1}{2} & -\frac{1}{2} & \frac{1}{2} \\ 
				\frac{1}{2} & \frac{1}{2} & -\frac{1}{2} \\ 
				-\frac{1}{2} & \frac{1}{2} & \frac{1}{2} 
			\end{pmatrix}
		\end{equation}
		To test this out, let us take the vector 
		\begin{equation}
			v = 1 \begin{pmatrix} 1 \\ 0 \\ 0 \end{pmatrix} + 1 \begin{pmatrix} 0 \\ 1 \\ 0 \end{pmatrix} + 1 \begin{pmatrix} 0 \\ 0 \\ 1 \end{pmatrix} = \frac{1}{2} \begin{pmatrix} 1 \\ 1 \\ 0 \end{pmatrix} + \frac{1}{2} \begin{pmatrix} 0 \\ 1 \\ 1 \end{pmatrix} + \frac{1}{2} \begin{pmatrix} 1 \\ 0 \\ 1 \end{pmatrix} 
		\end{equation}
		So $(1, 1, 1)^T$ should map to $(\frac{1}{2}, \frac{1}{2}, \frac{1}{2})^T$. Indeed, it is the case that 
		\begin{equation}
		  \begin{pmatrix}
				\frac{1}{2} \\ \frac{1}{2} \\ \frac{1}{2}
		  \end{pmatrix} = \begin{pmatrix}
				\frac{1}{2} & -\frac{1}{2} & \frac{1}{2} \\ 
				\frac{1}{2} & \frac{1}{2} & -\frac{1}{2} \\ 
				-\frac{1}{2} & \frac{1}{2} & \frac{1}{2} 
			\end{pmatrix} \begin{pmatrix} 
				1 \\ 1 \\ 1
			\end{pmatrix}
		\end{equation}
	\end{example}

	Given $n$-dimensional vector space $W$ over $\mathbb{F}$, let's say that we have linear map $A: W \to W$ and vector $w \in W$, with $w^\prime = Aw$. Given basis $\{u_1, \ldots, u_n\}$, we can write the vector $w_v \in \mathbb{R}^n$ with respect to this basis $\{u_i\}$. Furthermore, if we endow this basis to \textit{both} the domain $W$ and codomain $W$ of $A$, then we can write $A_{uu} \in \mathbb{F}^{n \times n}$ (where the left index represents the codomain and the right index represents the domain). This allows us to correctly compute 
	\begin{equation}
		w^\prime_u = A_{uu} w_u
	\end{equation}

	since they are all with respect to the same basis. If we have a second basis $\{v_1, \ldots, v_n\}$ of $V$ and choose to represent $w$ with respect to this basis, we will have $w_v = S w_u$, where $S$ is the change of basis matrix. But we can't just replace $w_u$ with $w_v$ since they are not with respect to the same basis. This won't work. 
	\begin{equation}
		w^\prime_u \neq A_{uu} w_v
	\end{equation} 
	We must change the basis of the domain of $A$ to $\{v_i\}$ in order for this to work. So we can do 
	\begin{align}
		w_u^\prime & = (A_{uu} S^{-1}) (S w_u) \\ 
							 & = A_{uv} s_v
	\end{align} 
	Furthermore, if we wanted to change the basis of the codomain, i.e. change the representation of $w^\prime$, then we can just multiply $S$ to $w_u^\prime$. 
	\begin{align}
		w_v^\prime = S w_u^\prime & = (S A_{uu} S^{-1}) (S w_u) \\ 
															& = A_{vv} s_v
	\end{align}

	This new matrix $A_{vv} = S A_{uu} S^{-1}$ is simply another representation of the same linear map $A$ with respect to a different basis (in both the domain and codomain). To repeat again, $A_{vv}$ and $A_{ww}$ represent the same transformation $A$ but merely in different bases. 

  \begin{definition}[Similar Matrices]
    Two matrices $A$ and $B$ are \textbf{similar} if and only if there exists an invertible matrix $S$ such that $B = S A S^{-1}$. 
  \end{definition}

	\begin{lemma}[Similarity is an Equivalence Relation]
		Matrix similarity is an \hyperref[st-def:equivalence_relation]{equivalence relation} on $\mathbb{F}^{n \times n}$. 
	\end{lemma}
	\begin{proof}
	  We prove the three properties of an equivalence relation. 
		\begin{enumerate}
			\item \textit{Reflexive}. $A \sim A$ since $A = I A I^{-1}$. 
			\item \textit{Symmetric}. If $A \sim B$, then $B = S A S^{-1}$, so $A = S^{-1} B S = S^{-1} B (S^{-1})^{-1}$, and so $B \sim A$. 
			\item \textit{Transitive}. If $A \sim B, B \sim C$, then $B = SAS^{-1}$ and $C = R B R^{-1}$, so $C = R S A S^{-1} R^{-1} = (RS) A (RS)^{-1} \implies A \sim C$. 
		\end{enumerate}
	\end{proof}

\subsection{Transpose and Adjoint Operators} 

  \begin{definition}[Transpose of a Matrix]
    The \textbf{transpose} of matrix $A$, denoted $A^T$, is the matrix with entries $(a^T)_{i j} = a_{i j}$. That is, it is $A$, "flipped over" its diagonal. 
  \end{definition}

  We illustrate why this definition of a transpose is equivalent to the abstract definition to the transpose of a linear map. Given a linear map $A: U \longrightarrow V$ with $\dim U = n, \dim V = m$, we can fix a basis on both $U$ and $V$ to define its matrix $A$. The abstract definition states that 
  \begin{equation}
    A^T: V^\ast \longrightarrow U^\ast, \; l \equiv \varphi A
  \end{equation}
  Treating $l$ and $\varphi$ as row vectors, we can see that the $m \times n$ matrix $A$ maps the $1 \times m$ covector $\varphi$ to the $1 \times n$ covector $l$. Note that this linear mapping is realized through \textit{right multiplication} of $A$ on $\varphi$. It is customary to present linear maps as \textit{left} multiplication, so by "flipping" (i.e. taking the matrix transpose) of all the elements in the equation, we get 
  \begin{equation}
    l^T \equiv A^T \varphi^T
  \end{equation}
  which presents the mapping in the more usual way of left matrix multiplication. Note that $l^T$ and $\varphi^T$ are still covectors. Just because they are now represented as column vectors, it does not mean that they are not covectors, which is why we shouldn't be too dependent on the row vector interpretation of dual vectors mentioned above.  

  \begin{theorem}[Properties of the Matrix Transpose]
    Given that $A, B: U \longrightarrow V$ is linear, $c$ a constant
    \begin{enumerate}
      \item $(A^T)^T = A$. 
      \item $(A+B)^T = A^T + B^T, \; (c A)^T = c A^T$. 
      \item $(A B)^T = B^T A^T$. 
      \item If $A$ is invertible, $(A^{-1})^T = (A^T)^{-1}$ and $A$ invertible $\implies A^T$ invertible. 
      \item $x \cdot y = x^T y$. Furthermore, 
      \begin{equation}
        Ax \cdot y = (A x)^T y = x^T A^T y = x \cdot A^T y
      \end{equation}
    \end{enumerate}
  \end{theorem}

	\begin{definition}[Symmetric Matrices]
    Matrix $A$ is 
		\begin{enumerate}
		  \item a \textbf{symmetric matrix} if $A = A^T$. 
			\item a \textbf{skew-symmetric} or \textbf{anti-symmetric} matrix if $A^T = - A$. 
		\end{enumerate}
  \end{definition}
	
  \begin{definition}[Row Space]
    The \textbf{row space} of matrix $A$, denoted $R(A)$, is the span of its row vectors. That is, 
    \begin{equation}
      R(A) = \text{span}\{ A_1, A_2, \cdots, A_m\}
    \end{equation}
    We will denote the row vectors with uppercase $A_i$'s. 
  \end{definition}

	\begin{definition}[Left Nullspace]
    The \textbf{left nullspace} of matrix $A$ is the nullspace of $A^T$. It is denoted $N(A^T)$. 
  \end{definition}

  \begin{theorem}
    Let $A: \mathbb{F}^n \longrightarrow \mathbb{F}^m$ be a $m \times n$ matrix with rank $k$. Assuming that $\mathbb{F}^n$ and $\mathbb{F}^m$ are inner product spaces,  
    \begin{align}
      N(A) = R(A)^\perp &\iff N(A)^\perp = R(A) \\
      N(A^T) = C(A)^\perp &\iff N(A^T)^\perp = C(A)
    \end{align}
    That is, $N(A)$ and $R(A)$ are orthogonal complements in $\mathbb{F}^n$, with $\dim R(A) = k$ and $\dim(N(A)) = n - k$. $N(A^T)$ and $C(A)$ are orthogonal complements in $\mathbb{F}^m$, with $\dim C(A) = k$ and $\dim(N(A^T)) = m - k$. 
  \end{theorem}

	\begin{theorem}[Column and Row Spaces are Isomorphic]
		Given matrix 
    \begin{equation}
      C(A) \simeq R(A)
    \end{equation}
  \end{theorem}
	\begin{proof}
	  The dimensions of the two subspaces are equal to the rank, and so they are isomorphic. 
	\end{proof}
  \begin{proof}
		We will prove that the linear mapping $A$ is the isomorphism itself. Let $\rank(A) = r$ and let $\{ v_1, v_2, \cdots, v_r\}$ be a basis for $R(A)$. Then, the set $\{ A v_1, A v_2, \cdots, A v_r\}$ are $r$ vectors in $C(A)$. They are linearly independent because 
    \begin{align*}
      \sum_{i = 1}^r c_i A v_i = A \sum_{i=1}^r c_i v_i = 0 & \implies \sum_{i=1}^r c_i v_i \in N(A) \text{, but } \sum_{i=1}^r c_i v_i \in R(A) \\
      & \implies \sum_{i=1}^r \in N(A) \cap R(A) = \{0\}
    \end{align*} 
    Since $\dim C(A) = r$, $\{A v_i\}$ must form a basis of $C(A)$. Therefore, $A$ is a bijection between vector spaces and is thus an isomorphism. 
  \end{proof}

  \begin{corollary}
    \begin{equation}
      \rank(A) = \rank(A^T)
    \end{equation}
  \end{corollary}

\subsection{Decomposition of Matrices}

\subsubsection{LUP Decomposition}

  \begin{theorem}
    The product of square lower triangular matrices is a lower triangular matrix. The product of square upper triangular matrices is an upper triangular matrix. 
  \end{theorem}

  \begin{theorem}[LU Decompositions]
		Every $A \in \mathbb{F}^{m \times n}$ admits a factorization, called the \textbf{LU decomposition}, in the form 
		\begin{equation}
		  A = LU
		\end{equation}
		where $L$ is a lower triangular matrix and $U$ is an upper triangular matrix. 
  \end{theorem}
  \begin{proof}
    We reduce $A$ to its echelon form ref$(A)$ by successively multiplying elementary matrices $E^{\gamma_i}$ representing elementary operation (i). After a finite amount of steps $r$, we will reduce it to ref$(A)$.
    \begin{equation}
      \text{ref}(A) = E^{\gamma_r} E^{\gamma_{r-1}} \cdots E^{\gamma_2} E^{\gamma_1} A = \bigg(\prod_{i = 0}^{r-1} E^{\gamma_{r-i}}\bigg) A
    \end{equation}
    Since each $E^{\gamma_i}$ is invertible, we multiply the product of the inverses of the elementary matrices of operation (i), which are also elementary matrices of operation (i). 
    \begin{equation}
      (E^{\gamma_1})^{-1} (E^{\gamma_2})^{-1} \cdots (E^{\gamma_r})^{-1} ref(A) = \bigg( \prod_{j = 1}^r (E^{\gamma_j})^{-1} \bigg) \bigg( \prod_{i=1}^{r-1} E^{\gamma_{r-i}} \bigg) A = A
    \end{equation}
    Since each $(E^{\gamma_j})^{-1}$ is an elementary row operation, it is lower diagonal, and by theorem 3.16, their product is also lower triangular. It is easy to prove that if the diagonal entries are furthermore equal to 1, then the product has diagonal entries equal to 1. Finally, it is clear that every matrix in row echelon form is upper triangular, and we are done. 
    \begin{equation}
      A = \bigg( \prod_{j = 1}^r (E^{\gamma_j})^{-1} \bigg) \text{ref}(A) = L U
    \end{equation}
  \end{proof}

  Note that the existence of the LU decomposition for a general $m \times n$ matrix is not guaranteed. It will not exist if we must switch rows in matrix $A$ in order to reduce it to its echelon form. It does not matter whether we need to use elementary operation (ii) or not. Only the necessity of elementary operation (iii) to reduce the matrix determines the existence of the LU decomposition. The decomposition is also unique. 

  Finding the LU decomposition of a matrix is useful for solving systems of linear equations. Given a system in the form of $A x = b$, if we know the LU decomposition of $A$, we can rewrite the system as 
  \begin{equation}
    L U x = b
  \end{equation}
  Setting $y = U x$, we can easily solve the system $L y = b$ using forward substitution and then we can solve the system $U x = y$ using back substitution. Therefore, knowing this decomposition beforehand greatly aids in computing the solutions to the linear system. But computing $L$ and $U$ in order to solve this system takes as much effort as solving the system using Gauss Elimination in the first place. 

  It is imperative to mention a similar decomposition for $n \times n$ matrices, known as \textbf{LUP decomposition}. 

  \begin{definition}
    An $n \times n$ permutation matrix is a matrix of $0$s and $1$s with exactly one $1$ in each row and column. The set of all $n \times n$ permutation matrices form a multiplicative matrix group of order $n!$. We can also view this group as the matrix representation of the symmetric group $S_n$. 
  \end{definition}

  \begin{example}
    The set of all $2 \times 2$ permutation matrices is 
    \begin{equation}
      S_2 = \bigg\{  \begin{pmatrix}1 & 0 \\0 & 1 \end{pmatrix}, \begin{pmatrix}1 & 0 \\0 & 1 \end{pmatrix} \bigg\}
    \end{equation}
    and the set of all $3 \times 3$ permutation matrices is
    \begin{equation}
      S_3 = \Bigg\{I_3, 
      \begin{pmatrix}1 & 0 & 0 \\0 & 0 & 1 \\0 & 1 & 0 \end{pmatrix},
      \begin{pmatrix}0 & 1 & 0 \\1 & 0 & 0 \\0 & 0 & 1 \end{pmatrix}, 
      \begin{pmatrix}0 & 1 & 0 \\0 & 0 & 1 \\1 & 0 & 0 \end{pmatrix}, 
      \begin{pmatrix}0 & 0 & 1 \\0 & 1 & 0 \\1 & 0 & 0 \end{pmatrix}, 
      \begin{pmatrix}0 & 0 & 1 \\1 & 0 & 0 \\0 & 1 & 0 \end{pmatrix} \Bigg\}
    \end{equation}
  \end{example}

  \begin{theorem}
    Every $n \times n$ matrix $A$ can be decomposed into the form $A = P L U$, where $L$ is lower triangular, $U$ is upper triangular, and $P$ is a permutation matrix. 
  \end{theorem}
  \begin{proof}
    We can modify the Gauss Elimination algorithm to do all the row interchanges in the beginning. The permutation matrices form a group, so the product of all the initial row changes is a permutation matrix. Call it $P^\prime$. The previous theorem states that we can do LU decomposition on $P^\prime A$. 
    \begin{equation}
      P^\prime A = LU \implies A = P^{\prime -1} L U = P L U
    \end{equation}
    Since $P^{\prime -1}$ is also in the symmetric group of permutations, we can denote it as $P$. 
  \end{proof}

  \begin{corollary}
    Every $n \times n$ matrix $A$ can be decomposed into the form $L U P$. That is, in the form
    \begin{equation}
      A = L U P = 
      \begin{pmatrix}
      1 & 0 & 0 & \ldots & 0\\
      * & 1 & 0 & \ldots & 0\\
      * & * & 1 & \ldots & \vdots\\
      \vdots & \vdots & \vdots & \ddots & 0\\
      * & * & \cdots & * & 1 
      \end{pmatrix}
      \begin{pmatrix}
      u_{11} & * & * & \ldots & *\\
      0 & u_{22} & * & \ldots & *\\
      0 & 0 & u_{33} & \ldots & \vdots\\
      \vdots & \vdots & \vdots & \ddots & * \\
      0 & 0 & \ldots & 0 & u_{n n} 
      \end{pmatrix}
      \begin{pmatrix}
      \\
      \\
       & & & P & & & \\
      \\
       & 
      \end{pmatrix}
    \end{equation}
  \end{corollary}
  \begin{proof}
    We decompose $A^T = P_0 L_0 U_0$, where $P_0$ is a permutation matrix, $L_0$ lower triangular, $U_0$ upper triangular. This implies that 
    \begin{equation}
      A = A^{T T} = U_0^T L_0^T P_0^T = L U P
    \end{equation}
    since $U_0^T$ is lower triangular and $L_0^T$ is upper triangular. Note that $L$ is unique, but $U$ is not unique, so this decomposition is not unique. 
  \end{proof}

  This decomposition can also be used to solve matrix equations
  \begin{equation}
    A X = B
  \end{equation}
  Since this equation can be expressed in the form
  \begin{equation}
    A \begin{pmatrix}
    | & & | \\ x_1 & \cdots & x_n \\ | & & | 
    \end{pmatrix} = \begin{pmatrix}
    | & & | \\ A x_1 & \cdots & A x_n \\ | & & | 
    \end{pmatrix} = \begin{pmatrix}
    | & & | \\ b_1 & \cdots & b_n \\ | & & | 
    \end{pmatrix}
  \end{equation}
  solving this matrix is equivalent to solving the system of systems of linear equations 
  \begin{equation}
    Ax_1 = b_1, Ax_2 = b_2, \cdots, Ax_n = b_n
  \end{equation}
  i.e. by solving one column at a time. This method can also be used to solve 
  \begin{equation}
    A X = I
  \end{equation}
  to find $X = A^{-1}$. Equivalently, we can left multiply elementary matrices to reduce $A$ to rref$(A)$. 
  \begin{equation}
    E^{\gamma_r} E^{\gamma_{r-1}} \cdots E^{\gamma_2} E^{\gamma_1} A X = \text{rref}(A) X = E^{\gamma_r} E^{\gamma_{r-1}} \cdots E^{\gamma_2} E^{\gamma_1} I = \prod_{i = 0}^{r-1} E^{\gamma_{r-i}}
  \end{equation}
  If rref$(A)= I$, then 
  \begin{equation}
    A^{-1} = \prod_{i = 0}^{r-1} E^{\gamma_{r-i}}
  \end{equation}
  and if rref$(A) \neq I$, then $A^{-1}$ does not exist. This is in fact precisely the method of finding the inverse where we do Gauss Elimination on the extended matrix
  \begin{equation}
    \begin{pmatrix}
    & &| & & \\ & A &| & I & \\ & &| & & 
    \end{pmatrix} \longrightarrow \begin{pmatrix}
    & &| & & \\ & I &| & A^{-1} & \\ & &| & & 
    \end{pmatrix}
  \end{equation}

\subsubsection{QR Decomposition} 

  The QR decomposition is often used to simplify these linear least squares problems into a more manageable equation. 

  \begin{definition}[QR Decomposition]
    Any $m \times n$ matrix $A$ over $\mathbb{F}$ (where $m \geq n$) may be decomposed into the product
    \begin{equation}
      A = QR
    \end{equation}
    where $Q$ is an $m \times n$ matrix with pairwise orthonormal column vectors (meaning $Q^\dagger Q = I_n$), and $R$ is an $n \times n$ upper triangular matrix. If $A$ has full column rank ($\rank(A) = n$), the decomposition is unique provided that we require the diagonal entries of $R$ to be strictly positive. 
  \end{definition}
  \begin{proof}
    We use the Gram-Schmidt orthogonalization process on the column vectors of $A$. Let $a_1, a_2, \dots, a_n$ be the columns of $A$. We can recursively construct the orthonormal columns $q_1, \dots, q_n$ of $Q$. 
    
    For the first step, let $v_1 = a_1$. Then $q_1 = v_1 / \|v_1\|$. We set $r_{11} = \|v_1\|$, giving $a_1 = r_{11} q_1$. 

    For step $j$ (where $1 < j \leq n$), we subtract the projections of $a_j$ onto the previously computed orthogonal vectors $q_1, \dots, q_{j-1}$:
    \begin{equation}
      v_j = a_j - \sum_{i=1}^{j-1} (a_j, q_i) q_i
    \end{equation}
    If $A$ has full column rank, then $a_j$ is not in the span of the previous vectors, so $v_j \neq 0$. We define $q_j = v_j / \|v_j\|$. 

    Let $r_{ij} = (a_j, q_i)$ for $i < j$ and $r_{jj} = \|v_j\|$. This allows us to write $a_j$ as:
    \begin{equation}
      a_j = \sum_{i=1}^j r_{ij} q_i
    \end{equation}
    This is precisely the matrix multiplication $A = QR$, where $R$ is upper triangular and $Q$ has orthonormal columns $q_i$. The condition $r_{jj} = \|v_j\| > 0$ makes the decomposition unique for full-rank matrices.
  \end{proof}

  Because $Q^\dagger Q = I$, the QR decomposition can greatly simplify normal equations in linear least squares problems:
  \begin{align}
    A^T A x = A^T b & \implies (Q R)^T (Q R) x = R^T Q^T Q R x = R^T R x = R^T Q^T b \\
    & \implies R x = Q^T b \\
    & \implies x = R^{-1} Q^T b
  \end{align}

  The classical Gram-Schmidt process is often numerically unstable. 

  \begin{algo}[Modified Gram-Schmidt for QR Decomposition]
    The Modified Gram-Schmidt algorithm provides better numerical stability by orthogonalizing the remaining column vectors against the newly computed orthonormal vector immediately at each step. 
    \begin{algorithmic}[1]
      \FOR{$j = 1, \dots, n$}
        \STATE $v_j = a_j$
      \ENDFOR
      \FOR{$i = 1, \dots, n$}
        \STATE $r_{ii} = \|v_i\|$
        \STATE $q_i = v_i / r_{ii}$
        \FOR{$j = i+1, \dots, n$}
          \STATE $r_{ij} = (v_j, q_i)$
          \STATE $v_j = v_j - r_{ij} q_i$
        \ENDFOR
      \ENDFOR
    \end{algorithmic}
  \end{algo}

\subsubsection{Cholesky Decomposition}

  \begin{definition}[Cholesky Decomposition]
    Every Hermitian (or real symmetric), positive-definite matrix $A$ can be decomposed into the product
    \begin{equation}
      A = L L^\dagger
    \end{equation}
    where $L$ is a lower triangular matrix with strictly positive real diagonal entries. This decomposition is unique.
  \end{definition}

  \begin{proof}
    We proceed by induction on the dimension $n$ of the matrix $A$. 
    
    \textbf{Base case:} For $n=1$, $A = (a_{11})$. Since $A$ is positive-definite, $a_{11} > 0$. We simply choose $L = (\sqrt{a_{11}})$, which is unique and positive. 
    
    \textbf{Inductive step:} Assume the theorem holds for $(n-1) \times (n-1)$ matrices. We partition the $n \times n$ matrix $A$ as:
    \begin{equation}
      A = \begin{pmatrix} a_{11} & w^\dagger \\ w & A_{22} \end{pmatrix}
    \end{equation}
    where $a_{11}$ is a positive scalar (since $A$ is positive-definite, its diagonal entries are positive), $w$ is an $(n-1) \times 1$ column vector, and $A_{22}$ is an $(n-1) \times (n-1)$ Hermitian matrix. 
    
    We seek a lower triangular matrix $L$ of the form:
    \begin{equation}
      L = \begin{pmatrix} l_{11} & 0 \\ l_{21} & L_{22} \end{pmatrix}
    \end{equation}
    such that $A = L L^\dagger$. Multiplying out $L L^\dagger$, we get:
    \begin{equation}
      \begin{pmatrix} l_{11}^2 & l_{11} l_{21}^\dagger \\ l_{11} l_{21} & l_{21} l_{21}^\dagger + L_{22} L_{22}^\dagger \end{pmatrix} = \begin{pmatrix} a_{11} & w^\dagger \\ w & A_{22} \end{pmatrix}
    \end{equation}
    Matching the blocks, we uniquely determine $l_{11} = \sqrt{a_{11}}$ (since $l_{11} > 0$). Then we can uniquely determine the column vector $l_{21} = w / l_{11}$. 
    
    Finally, we must satisfy $L_{22} L_{22}^\dagger = A_{22} - l_{21} l_{21}^\dagger$. The matrix $S = A_{22} - l_{21} l_{21}^\dagger$ is the Schur complement of $a_{11}$ in $A$. Because $A$ is positive-definite, its Schur complement $S$ is also positive-definite. By our inductive hypothesis, the $(n-1) \times (n-1)$ positive-definite matrix $S$ has a unique Cholesky decomposition $S = L_{22} L_{22}^\dagger$. Thus, $L_{22}$ exists and is unique, completing the proof.
  \end{proof}

  \begin{algo}[Cholesky-Banachiewicz Algorithm]
    Intuitively, this algorithm computes the lower triangular matrix $L$ starting from the upper left corner and proceeding row by row. It requires about half the operations of the LU decomposition.
    \begin{algorithmic}[1]
      \FOR{$i = 1, \dots, n$}
        \FOR{$j = 1, \dots, i$}
          \IF{$j = i$}
            \STATE $l_{jj} = \sqrt{ a_{jj} - \sum_{k=1}^{j-1} |l_{jk}|^2 }$
          \ELSE
            \STATE $l_{ij} = \frac{1}{l_{jj}} \left( a_{ij} - \sum_{k=1}^{j-1} l_{ik} \bar{l}_{jk} \right)$
          \ENDIF
        \ENDFOR
      \ENDFOR
    \end{algorithmic}
  \end{algo}

\subsection{Block Matrices}

  \begin{theorem}
    Given mappings $A_i \in$ End$(V_i)$ for $i = 1, 2, \cdots, n$, the matrix representation of the induced linear mapping $A_1 \oplus A_2 \oplus \cdots \oplus A_n$ is the block matrix 
    \begin{equation}
      \begin{pmatrix}
      A_1 & & & \\
      & A_2 & & \\
      & & \ddots & \\
      & & & A_n
      \end{pmatrix}: \bigoplus_{i=1}^n V_i \longrightarrow \bigoplus_{i=1}^n V_i
    \end{equation}
  \end{theorem}

  \begin{theorem}
    Given $2 \times 2$ block matrices
    \begin{equation}
      X = \begin{pmatrix}
      A_1&B_1\\C_1&D_1
      \end{pmatrix}, \; \; Y = \begin{pmatrix}
      A_2&B_2\\C_2&D_2
      \end{pmatrix}
    \end{equation}
    We can compute $X Y$ similarly to regular matrix multiplication, treating the blocks as entries. 
    \begin{equation}
      X Y = \begin{pmatrix}
      A_1&B_1\\C_1&D_1
      \end{pmatrix} \begin{pmatrix}
      A_2&B_2\\C_2&D_2
      \end{pmatrix} = \begin{pmatrix}
      A_1 A_2 + B_1 C_2 & A_1 B_2 + B_1 D_2 \\
      C_1 A_2 + D_1 C_2 & C_1 B_2 + D_1 D_2 
      \end{pmatrix}
    \end{equation}
    Furthermore, this process can be done in general for any $m \times n$ block matrix $X$ and $n \times p$ block matrix $Y$. 
  \end{theorem}

  \begin{theorem}
    Given that $I_N, A, B$ are $n \times n$ matrices, define the $(2n) \times (2n)$ matrix 
    \begin{equation}
      X = \begin{pmatrix} I & 0 \\ A & B \end{pmatrix}
    \end{equation}
    Then 
    \begin{equation}
      \det{X} = \det{B}
    \end{equation}
  \end{theorem}
  \begin{proof}
    We can perform Gauss elimination to reduce $X$ without affecting the determinant.
    \begin{equation}
      \det{\begin{pmatrix} I&0\\A&B \end{pmatrix}} = \det{ \begin{pmatrix} I&0\\ 0&B \end{pmatrix}} = \det{B}
    \end{equation}
    since it satisfies the correct properties for $\det{B}$. 
  \end{proof}

  \begin{corollary}
    \begin{equation}
      \det{\begin{pmatrix} A&0\\C&D \end{pmatrix}} = \det{A} \det{D}
    \end{equation}
  \end{corollary}
  \begin{proof}
    \begin{equation}
      \det{\begin{pmatrix}
      A&0\\C&D 
      \end{pmatrix}} = \det{\begin{pmatrix}
      A&0\\C&I
      \end{pmatrix} \begin{pmatrix}
      I&0\\0&D
      \end{pmatrix}} = \det{\begin{pmatrix}
      A&0\\C&I
      \end{pmatrix}} \det{\begin{pmatrix}
      I&0\\0&D
      \end{pmatrix}}
    \end{equation}
  \end{proof}

  However, 
  \begin{equation}
    \det{\begin{pmatrix} A&B\\C&D \end{pmatrix}} \neq \det{A} \det{D} - \det{B} \det{C}
  \end{equation}

  Rather, we introduce the following theorem

  \begin{theorem}
    \begin{align}
      \det{\begin{pmatrix} A&B\\C&D \end{pmatrix}}  & = \det{(A)} \det{(D - C A^{-1} B)} \\
      & = \det{(D)} \det{(A - B D^{-1} C)}
    \end{align}
  \end{theorem}
  \begin{proof}
    \begin{equation}
      \begin{pmatrix} A&B\\C&D\end{pmatrix} = \begin{pmatrix}
      A&0\\C&I\end{pmatrix} \begin{pmatrix}
      I& A^{-1} B \\ 0 & D - C A^{-1} B
      \end{pmatrix}
    \end{equation}
    by similarity, equation $(6)$ is equal to equation $(7)$. 
  \end{proof}

  \begin{definition}[Block Diagonal Matrix]
    A \textbf{block diagonal matrix} is a square matrix in block form such that the diagonal blocks are square matrices and all off-diagonal blocks are zero matrices. 
    \begin{equation}
      A = \begin{pmatrix}
      A_1&0&\ldots&0\\
      0&A_2&\ldots&0\\
      \vdots&\vdots&\ddots&\vdots\\
      0&0&\ldots&A_k
      \end{pmatrix}
    \end{equation}
  \end{definition}

  \begin{theorem}
    Given a matrix $A$ in block diagonal form, with diagonal blocks $A_1, A_2, \cdots, A_k$,
    \begin{equation}
      \det{A} = \prod_{i=1}^k A_i, \; \; \Tr{A} = \sum_{i=1}^k \Tr{A_i}
    \end{equation}
    Furthermore, $A$ is invertible if and only if all the $A_i$'s are invertible, and 
    \begin{equation}
      A^{-1} = \begin{pmatrix}
      A_1&0&\ldots&0\\
      0&A_2&\ldots&0\\
      \vdots&\vdots&\ddots&\vdots\\
      0&0&\ldots&A_k
      \end{pmatrix} = \begin{pmatrix}
      A_1^{-1}&0&\ldots&0\\
      0&A_2^{-1}&\ldots&0\\
      \vdots&\vdots&\ddots&\vdots\\
      0&0&\ldots&A_k^{-1}
      \end{pmatrix}
    \end{equation}
  \end{theorem}
  \begin{proof}
    The results are obvious when performing block multiplication or Gauss Elimination. 
  \end{proof}

\subsection{Matrix Groups}

  \begin{definition}[Orthogonal Matrix]
    An \textbf{orthogonal matrix} is the matrix representation of an element in $O(n)$. It is the real $n \times n$ matrix where all the column vectors are pairwise orthogonal and all have magnitude 1. 

		The matrix representation of this group is the set of real $n \times n$ matrices where the column vectors form an orthonormal basis. Note that the determinant of every element of $O(n)$ is $\pm 1$. 
  \end{definition}

  \begin{theorem}
    The rows of an orthogonal matrix are also pairwise orthonormal.
  \end{theorem}

	\begin{definition}[Special Orthogonal Matrix]
	  The matrix representation of this group is the set of real $n \times n$ matrices where the column vectors are pairwise orthonormal and the determinant $=1$. 
	\end{definition}

  \begin{theorem}
    Given an orthogonal matrix $M$,
    \begin{equation}
      M^T = M^{-1}
    \end{equation}
  \end{theorem}

\subsection{Duality Theorem}

  The duality theorem allows us to define linear programming and establish the equivalence of primal/dual optimization. 

  In this section we will denote vector inequalities as entry-wise inequalities.Recall that elements of a vector space $X$ can be interpreted as column vectors, and elements of the dual of the vector space $X^\ast$ can be interpreted as row vectors. Therefore, value of $\phi$ at $x$ is denoted
  \begin{equation}
    \phi (x) = \phi_1 x_1 + \phi_2 x_2 + \cdots + \phi_n x_n
  \end{equation}

  Furthermore, the dual of $X^\ast$ is $X$ itself, and given that $Y$ is a linear subspace of $X$, the annihilator of $Y^\perp$ is $Y$. 
  \begin{equation}
    X = X^{**}, \qquad Y = Y^{\perp\perp}
  \end{equation}

  Suppose $Y$ is defined as the linear space spanned by the $m$ vectors $y_1, y_2, \cdots, y_m$ in $X$. That is, $Y$ consists of all vectors $y$ of the form
  \begin{equation}
    y = \sum_{j=1}^m a_j y_j
  \end{equation}

  It is clear by linearity that $\phi$ belongs to $Y^\perp$ if and only if
  \begin{equation}
    \phi (y_j) = 0, \qquad j = 1, 2, \cdots, m
  \end{equation}

  That is, a vector $y$ can be written as a linear combination of $m$ given vectors $y_j$ if and only if every $\phi$ that annihilates the $m$ vectors a also annihilates $0$. Now, we state a theorem that allows us to check if a vector $y$ can be written as a \textit{nonnegative} linear combinations of the $y_j$s. 

  \begin{theorem}
    A vector $y$ can be written as a linear combination of given vectors $y_j$ with nonnegative coefficients if and only if every $\zeta \in X^\ast$ that satisfies 
    \begin{equation}
       \zeta (y_j) \geq 0, \;\; j = 1, 2, \cdots, m
    \end{equation}
    also satisfies 
    \begin{equation}
      \zeta (y) \geq 0
    \end{equation}
  \end{theorem}
  \begin{proof}
    The proof is not the easiest to construct rigorously, but it can be visualized easily. 
  \end{proof}

  \begin{corollary}
    Given a $n \times m$ matrix $Y$, a vector $y$ with $n$ components can be written in the form 
    \begin{equation}
      y = Y p, \;\; p \geq 0
    \end{equation}
    if and only if every row vector $\zeta$ that satisfies 
    \begin{equation}
      \zeta Y \geq 0
    \end{equation}
    also satisfies 
    \begin{equation}
      \zeta y \geq 0
    \end{equation}
  \end{corollary}

  \begin{theorem}
    Given an $n \times m$ matrix $Y$ and a column vector $y$ with $n$ components, the inequality 
    \begin{equation}
      y \geq Y p, \;\; p \geq 0
    \end{equation}
    is satisfied if and only if every $\zeta$ that satisfies
    \begin{equation}
      \zeta Y \geq 0, \;\; \zeta \geq 0
    \end{equation}
    also satisfies 
    \begin{equation}
      \zeta y \geq 0
    \end{equation}
  \end{theorem}

  \begin{theorem}[Duality Theorem]
    Let $Y$ be a given $n \times m$ matrix, $y$ a given column vector with $n$ components, and $\gamma$ a given row vector with $m$ components. Let 
    \begin{equation}
      S = \sup_p \{\gamma p\}
    \end{equation}
    for all column vectors $p$ with $m$ components satisfying $y \geq Y p, \; p \geq 0$. A well-defined such $p$ is called \textbf{supremum admissible}. Additionally, let 
    \begin{equation}
      s = \inf_\zeta \{ \zeta y\}
    \end{equation}
    for all row vectors $\zeta$ with $n$ components satisfying $\gamma \leq \zeta Y, \; \zeta \geq 0$. A well-defined such $\zeta$ is called \textbf{infimum admissible}. Given that admissible vectors $p$ and $\zeta$ exist, then $S$ and $s$ are finite and 
    \begin{equation}
      S = s
    \end{equation}
  \end{theorem}

